\section{重—轻核的线性重整子}\label{sec:yzliu-renormalon}

重—轻核 $H_{HL}$ 的相位来自类空连线上眼形传播子的切割。Liu 与 Su\cite{Liu:2023onm}发现这一相位带有线性红外重整子，其强度与极点质量的线性重整子相同。

\begin{proposition}[重—轻 Sudakov 核的线性红外重整子 \tagP]\label{prop:yzliu-linear-renormalon}
在大 $\beta_0$（泡链）近似中，$\overline{\mathrm{MS}}$ 单顶点重—轻核 $C^{-\mathrm{sgn}z}_{HL}$（\cref{prop:yzliu-hard-kernel}(b)；即文献\cite{Liu:2023onm}的 $H_{HL}$）的 Borel 变换（Borel 变量 $u=\beta_0^{(b)}t/2$，与本书的 $u$ 相同，\cref{def:power-borel}）在右半平面的第一个奇点位于 $u=\frac12$，且只出现在虚部： % VF-YZL-11
\begin{equation}\label{eq:yzliu-renormalon}
\begin{gathered}
\Imag B(u)\xrightarrow{u\to\frac12}\frac{\mu}{p^z}\frac{C_F\ee^{5/6}}{2\pi}\frac{\mathrm{sgn}(z)}{u-\frac12},\\
A_{\mathrm{asym}}=-2\frac{\mu}{p^z}\sum_nr_n\alpha_s^{n+1},\quad r_n=N_m(\beta_0^{(b)})^n\Gamma(n+1),\quad N_m=\frac{C_F\ee^{5/6}}{\pi}\simeq0.977,
\end{gathered}
\end{equation}
$N_m$ 即极点质量与 Wilson 线线性重整子的归一化；实部没有线性重整子，轻—轻 Sudakov 核也没有。它 (i) 在类空夸克双线性系数函数的阈值展开（坐标空间与矩空间）中，与一个显式的次领头幂次阈值软算符的紫外重整子相消；(ii) 在准光前波函数的 TMD 展开中，与含量纲 1 的 Wilson 线尖点算符 $\hat O=[n_z\infty,0]\int_{-\infty}^0\dd x^+[0,x^+]D_\perp^2[x^+,-\infty]$ 的次领头幂次软因子 $S^\pm_q$ 的紫外重整子相消。该重整子与线性紫外发散相联，而与线性快度发散无关；在泡链层次，准光前波函数的 TMD 展开不收敛（留数正负交替，并有指数小项）；Coulomb 规范中的夸克场重整化也有线性重整子。（本命题只涉及部分子态与算符。）
\uses{prop:yzliu-hard-kernel,def:power-borel,prop:power-renormalon-examples,prop:lamet-lrr}
\end{proposition}

\begin{proofsketch}
（1）泡链：把单圈图（\cref{prop:yzliu-hard-kernel}(b)）中的胶子传播子换成链和；链和的 Borel 变换为 $(\mu^2\ee^{5/3})^u(-k^2)^{-1-u}$\cite[式~(2.44)--(2.45)]{Liu:2023onm}。（2）虚部来自对眼形传播子作切割 $\frac1{k^z\pm\ii0}\to\mp\ii\pi\delta(k^z)$；余下的积分给出 $\Gamma$ 函数之比，在 $u=\frac12$ 有单极点，留数为 $\frac{C_F\,\mathrm{sgn}(z)}{\pi}\ee^{5/6-L_\zeta/2}=\frac{\mu}{p^z}\frac{C_F\ee^{5/6}}{2\pi}\mathrm{sgn}(z)$（$\ee^{-L_\zeta/2}=\mu/(2p^z)$），即\cref{eq:yzliu-renormalon}\cite[式~(2.46)--(2.47)]{Liu:2023onm}；实部在 $u=\frac12$ 正则。（3）由 Borel 变换的定义（\cref{def:power-borel}），单极点 $\frac{c}{u-\frac12}$ 对应系数 $-2c\,(\beta_0^{(b)})^kk!$；按 $\mathrm{sgn}(z)A=2\,\mathrm{Arg}\,C^{-\mathrm{sgn}z}_{HL}$ 换算即得 $A_{\mathrm{asym}}$ 与 $N_m$\cite[式~(2.48)--(2.50)]{Liu:2023onm}；歧义 $\sim(\mu/p^z)(\Lambda^2/\mu^2)^{1/2}=\Lambda/p^z$，是线性幂次修正。（4）相消：阈值展开见\cite[第~3.6~节]{Liu:2023onm}，准光前波函数见\cite[第~4.3--4.4~节、式~(4.28)--(4.33)]{Liu:2023onm}。启发式来源：类空连线产生的“虚势” $gA^z\sim\ii V$ 是线性静态势的解析延拓，给出 $-\ii\frac{\delta m}{2p^z}$ 型修正。
\end{proofsketch}

\paragraph{地位与缺口。}[P]：在泡链中精确，超出大 $\beta_0$ 无推导（\cref{conj:yzliu-renormalon-norm}）。推论：它解释了格点 CS 核提取（\cref{prop:lattice-cs-kernel}）中虚部收敛缓慢的现象；取 Borel 重求和的硬核可减轻这一问题。背景见\cite{Beneke1999,Braun:2018brg}；Wilson 线自能的线性重整子见\cref{prop:lamet-lrr}。

\begin{conjecture}[线性重整子归一化的普适性 \tagC]\label{conj:yzliu-renormalon-norm}
在大 $\beta_0$ 之外，重—轻核相位的线性重整子归一化仍等于极点质量的 $N_m$：$n_f=0,\dots,5$ 时 $N_m=0.622,\dots,0.524$\cite[式~(2.52)]{Liu:2023onm}。（本猜想只涉及部分子态。）
\uses{prop:yzliu-linear-renormalon}
\end{conjecture}

\paragraph{地位与缺口。}[C]：文献\cite{Liu:2023onm}类比极点质量与 Wilson 线的线性重整子提出；只有泡链值被推导（\cref{prop:yzliu-linear-renormalon}），对重—轻核没有超出大 $\beta_0$ 的论证。升级途径：对重—轻核作超出泡链的重整子分析。

\begin{conjecture}[强子夸克双线性矩阵元的阈值渐近 \tagC]\label{conj:yzliu-threshold-hadronic}
在中等 $z$ 处，强子夸克双线性矩阵元在阈值极限 $\lambda\to-\ii\infty$（本书的 Ioffe 时间；Liu 的变量中为 $\lambda\to+\ii\infty$）有完全因子化的渐近形式：重—轻核 $H_{HL}$ 乘以非微扰阈值软因子 $J(z\Lambda_{\mathrm{QCD}})$ 与扭度 2 PDF 的阈值极限，修正为 $\order{1/|\lambda|}$\cite[式~(3.50)]{Liu:2023onm}。非微扰地位：\cref{asm:intro-wightman}、\cref{asm:intro-composite}（MA0）。
\uses{prop:yzliu-linear-renormalon,thm:yzliu-threshold}
\end{conjecture}

\paragraph{地位与缺口。}[C]：文献\cite{Liu:2023onm}的假设形式。理由是：欧氏时间演化的指数压低迫使大能量在两个夸克—连线顶点处流入流出，中间态（“割”）保持轻，故大 $|\lambda|$ 由阈值区主导。它预设了\cref{thm:yzliu-threshold}（只对部分子核为 [A]）的强子层次版本，以及非微扰阈值软因子 $J(z\Lambda_{\mathrm{QCD}})$；强子层次没有论证。升级途径：强子层次的阈值因子化（参见\cref{asm:intro-bridge}）。
